\documentclass[twocolumn]{aastex631}
\begin{document}
\title{A residual reionization ionizing-photon-budget shortfall for star-forming galaxies at z\~{}8 (lowxi systematic corner)}
\author{NebulaMind Lab (autonomous pipeline)}
\affiliation{NebulaMind Lab, \url{https://lab.nebulamind.net}}
\begin{abstract}
Reionization ionizing-photon-budget reconciliation at z\~{}8: star-forming galaxies require f\_esc=0.576 (+0.580/-0.295) to close the budget (Madau-Dickinson SFRD, log xi\_ion=25.5+/-0.15, clumping C=2-5, JWST-SFRD tail), versus indirect-proxy-inferred f\_esc=0.062 (+0.108/-0.039) from LzLCS O32/beta calibrations. Median delta(required-inferred)=+0.483 dex-frac (16-84\%: +0.176 to +1.064); 96\% of the systematic MC shows a shortfall. Conclusion: a genuine SHORTFALL remains, and the sign holds under both O32 and beta calibrations. The result is bounded by the xi\_ion x clumping x proxy-calibration systematic, not by statistics. Generated autonomously from public data (jwst) via the NebulaMind Lab runner: a bounded, reproducible, descriptive study.
\end{abstract}
\section{Introduction}
Introduction: Recent studies have highlighted a potential crisis in understanding the reionization process, with concerns raised about the sufficiency of ionizing photons produced by star-forming galaxies to drive this epochal transition [Muoz2024]. This issue is compounded by the need for increased demands on ionizing sources due to the absorption-dominated nature of reionization [Davies2021]. To address these challenges, it is essential to reconcile the ionizing photon budget using established literature values and analytic models. Our work builds upon previous efforts to calibrate excursion set reionization models [Park2022] and assesses the galaxy ionizing photon budget at high redshifts [Duncan2015].
\section{Data and method}
Data and method: We adopt a literature-anchored approach, utilizing the cosmic star formation rate density (SFRD) from Madau \& Dickinson's (2014) analytic fitting function. The ionization efficiency (xi\_ion) is set to log xi\_ion = 25.5 ± 0.15, and clumping factor (C) ranges between 2-5. To estimate the escape fraction (f\_esc), we rely on published calibrations from LzLCS O32/beta proxies [Chisholm+22, Flury+22; Simmonds+24]. Our method focuses on reconciling these parameters to determine if star-forming galaxies can account for the required ionizing photons during reionization.
\section{Result}
Result: Through our analysis, we find that star-forming galaxies require an escape fraction of f\_esc = 0.576 (+0.580/-0.295) to close the ionizing photon budget at z\~{}8. However, indirect-proxy-inferred values from LzLCS O32/beta calibrations yield a significantly lower f\_esc = 0.062 (+0.108/-0.039). This discrepancy results in a median shortfall of +0.483 dex-frac (16-84\%: +0.176 to +1.064), with 96\% of our systematic Monte Carlo simulations indicating a genuine shortfall. Notably, this result holds under both O32 and beta calibrations.
\begin{figure}\centering\includegraphics[width=\columnwidth]{result.png}\caption{A residual reionization ionizing-photon-budget shortfall for star-forming galaxies at z\~{}8 (lowxi systematic corner)}\end{figure}
\section{Caveats}
Caveats: Our study relies on an automated, single-selection, uncalibrated measurement approach, which has inherent limitations. Firstly, the use of published literature values introduces potential biases from differing methodologies and assumptions in the original studies. Secondly, our reliance on indirect-proxy-inferred escape fractions may not fully capture the complexities of ionizing photon escape in star-forming galaxies. Lastly, the clumping factor's uncertainty can significantly impact the required escape fraction, highlighting the need for further research to better constrain this parameter. These limitations underscore the importance of future observational and theoretical efforts to refine our understanding of reionization mechanisms.
\section*{References}
\footnotesize
[Muoz2024] Muñoz et al. (2024). Reionization after JWST: a photon budget crisis?. bibcode 2024MNRAS.535L..37M \par [Davies2021] Davies et al. (2021). The Predicament of Absorption-dominated Reionization: Increased Demands on Ionizing Sources. bibcode 2021ApJ...918L..35D \par [Park2022] Park et al. (2022). Calibrating excursion set reionization models to approximately conserve ionizing photons. bibcode 2022MNRAS.517..192P \par [Duncan2015] Duncan et al. (2015). Powering reionization: assessing the galaxy ionizing photon budget at z < 10. bibcode 2015MNRAS.451.2030D \par [Madau2017] Madau (2017). Cosmic Reionization after Planck and before JWST: An Analytic Approach. bibcode 2017ApJ...851...50M

\end{document}
